\begin{problem}[Conditional selection]
From a standard deck, one card is selected. Given that the card is red, find the probability that it is a queen.
\end{problem}

\begin{hint}
Restrict the sample space to the red cards first.
\end{hint}

\begin{solution}
There are $26$ red cards and $2$ red queens.
\[
P(\text{queen}\mid \text{red})=\frac2{26}=\frac1{13}.
\]
\end{solution}

\begin{problem}[Independence check]
A fair die is rolled once. Let
\[
A=\{\text{an even number}\},\qquad B=\{\text{a number greater than }3\}.
\]
Determine whether $A$ and $B$ are independent.
\end{problem}

\begin{hint}
Compute $P(A)$, $P(B)$, and $P(A\cap B)$, then compare.
\end{hint}

\begin{solution}
\[
P(A)=\frac36=\frac12,\qquad P(B)=\frac36=\frac12.
\]
Also,
\[
A\cap B=\{4,6\},
\]
so
\[
P(A\cap B)=\frac26=\frac13.
\]
Since
\[
\frac13\ne \frac12\cdot\frac12=\frac14,
\]
the events are not independent.
\end{solution}

\begin{problem}[Finding a missing probability]
A random variable $X$ has
\[
P(X=0)=0.2,\qquad P(X=1)=0.35,\qquad P(X=2)=p,\qquad P(X=3)=0.15.
\]
Find $p$ and then find $E(X)$.
\end{problem}

\begin{hint}
Probabilities must add to $1$ before you calculate expectation.
\end{hint}

\begin{solution}
\[
0.2+0.35+p+0.15=1
\]
gives
\[
p=0.30.
\]
Then
\[
E(X)=0(0.2)+1(0.35)+2(0.30)+3(0.15)=0.35+0.60+0.45=1.40.
\]
\end{solution}

\begin{problem}[Binomial mean and variance]
If
\[
X\sim \mathrm{Bin}(12,0.25),
\]
find $E(X)$ and $\mathrm{Var}(X)$.
\end{problem}

\begin{hint}
Use the standard formulas for a binomial random variable.
\end{hint}

\begin{solution}
\[
E(X)=np=12(0.25)=3.
\]
\[
\mathrm{Var}(X)=np(1-p)=12(0.25)(0.75)=2.25.
\]
\end{solution}

\begin{problem}[A union from a survey table]
In a group of $40$ students, $24$ study French, $18$ study Music, and $10$ study both. One student is selected at random. Find the probability that the student studies French or Music.
\end{problem}

\begin{hint}
Use the addition rule with counts.
\end{hint}

\begin{solution}
\[
P(F\cup M)=\frac{24+18-10}{40}=\frac{32}{40}=\frac45.
\]
\end{solution}

\begin{problem}[Variance from a probability table]
A random variable $X$ has
\[
P(X=0)=0.5,\qquad P(X=2)=0.3,\qquad P(X=4)=0.2.
\]
Find $\mathrm{Var}(X)$.
\end{problem}

\begin{hint}
First find $E(X)$ and $E(X^2)$.
\end{hint}

\begin{solution}
\[
E(X)=0(0.5)+2(0.3)+4(0.2)=1.4.
\]
\[
E(X^2)=0^2(0.5)+2^2(0.3)+4^2(0.2)=4.4.
\]
So
\[
\mathrm{Var}(X)=4.4-(1.4)^2=4.4-1.96=2.44.
\]
\end{solution}

\begin{problem}[Distribution table from draws without replacement]
A bag contains six balls marked
\[
0,1,2,3,4,5.
\]
Three balls are drawn in succession, without replacement, and the number of odd-numbered balls drawn is recorded. Let $X$ be the number of odd-numbered balls chosen.

Draw up a probability distribution table for $X$.
\end{problem}

\begin{hint}
There are $3$ odd-numbered balls and $3$ even-numbered balls. Use combinations:
\[
P(X=r)=\frac{\binom{3}{r}\binom{3}{3-r}}{\binom{6}{3}}.
\]
\end{hint}

\begin{solution}
The odd-numbered balls are
\[
1,3,5,
\]
and the even-numbered balls are
\[
0,2,4.
\]

Since $3$ balls are drawn, the possible values of $X$ are
\[
0,1,2,3.
\]

There are
\[
\binom63=20
\]
equally likely selections of $3$ balls.

\textbf{When $X=0$:} all $3$ chosen balls are even, so
\[
P(X=0)=\frac{\binom30\binom33}{\binom63}=\frac1{20}.
\]

\textbf{When $X=1$:} choose $1$ odd and $2$ even, so
\[
P(X=1)=\frac{\binom31\binom32}{\binom63}=\frac9{20}.
\]

\textbf{When $X=2$:} choose $2$ odd and $1$ even, so
\[
P(X=2)=\frac{\binom32\binom31}{\binom63}=\frac9{20}.
\]

\textbf{When $X=3$:} all $3$ chosen balls are odd, so
\[
P(X=3)=\frac{\binom33\binom30}{\binom63}=\frac1{20}.
\]

Therefore the probability distribution table is
\[
\begin{array}{c|cccc}
x & 0 & 1 & 2 & 3 \\
\hline
P(X=x) & \frac1{20} & \frac9{20} & \frac9{20} & \frac1{20}
\end{array}
\]
\end{solution}

\begin{problem}[A first binomial distribution table]
A fair coin is tossed $4$ times and the number of heads is recorded. Construct a table showing the probability distribution of $X$, the number of heads, and calculate $E(X)$.

\medskip
\noindent\textit{Note.} This is an example of a binomial distribution: a multi-stage experiment with identical stages, where at each stage there are two possible outcomes. Here there are $4$ stages because the coin is tossed $4$ times, and each stage has two possible outcomes, heads or tails.
\end{problem}

\begin{hint}
For a fair coin,
\[
X\sim \mathrm{Bin}\left(4,\frac12\right).
\]
Use
\[
P(X=r)=\binom{4}{r}\left(\frac12\right)^4.
\]
\end{hint}

\begin{solution}
Since the coin is fair and tossed $4$ times independently,
\[
X\sim \mathrm{Bin}\left(4,\frac12\right).
\]

So for $r=0,1,2,3,4$,
\[
P(X=r)=\binom{4}{r}\left(\frac12\right)^r\left(\frac12\right)^{4-r}
=\binom{4}{r}\left(\frac12\right)^4
=\frac{\binom{4}{r}}{16}.
\]

Therefore
\[
\begin{array}{c|ccccc}
x & 0 & 1 & 2 & 3 & 4 \\
\hline
P(X=x) & \frac1{16} & \frac4{16} & \frac6{16} & \frac4{16} & \frac1{16}
\end{array}
\]

Now
\[
E(X)=0\left(\frac1{16}\right)+1\left(\frac4{16}\right)+2\left(\frac6{16}\right)+3\left(\frac4{16}\right)+4\left(\frac1{16}\right).
\]
So
\[
E(X)=\frac{0+4+12+12+4}{16}=\frac{32}{16}=2.
\]

Hence the expected number of heads is
\[
\boxed{2}.
\]
\end{solution}

\begin{problem}[A first hypergeometric distribution table]
Two cards are drawn in succession from a standard deck, without replacement, and the number of hearts is recorded. Construct a table showing the probability distribution of $X$, the number of hearts drawn, and calculate $E(X)$.

\medskip
\noindent\textit{Note.} This is an example of a hypergeometric distribution: a multi-stage experiment with two possible outcomes at each stage, where selection is made without replacement. Here there are $2$ stages because two cards are drawn.
\end{problem}

\begin{hint}
The possible values are $0$, $1$, and $2$. Use combinations, or count the three cases directly.
\end{hint}

\begin{solution}
There are $13$ hearts and $39$ non-hearts in a standard deck of $52$ cards.

The possible values of $X$ are
\[
0,1,2.
\]

Using combinations,
\[
P(X=0)=\frac{\binom{39}{2}}{\binom{52}{2}}=\frac{741}{1326}=\frac{247}{442},
\]
\[
P(X=1)=\frac{\binom{13}{1}\binom{39}{1}}{\binom{52}{2}}=\frac{507}{1326}=\frac{169}{442},
\]
and
\[
P(X=2)=\frac{\binom{13}{2}}{\binom{52}{2}}=\frac{78}{1326}=\frac{13}{221}.
\]

Therefore the probability distribution table is
\[
\begin{array}{c|ccc}
x & 0 & 1 & 2 \\
\hline
P(X=x) & \frac{247}{442} & \frac{169}{442} & \frac{13}{221}
\end{array}
\]

Now
\[
E(X)=0\left(\frac{247}{442}\right)+1\left(\frac{169}{442}\right)+2\left(\frac{13}{221}\right).
\]
So
\[
E(X)=\frac{169}{442}+\frac{26}{221}
=\frac{169}{442}+\frac{52}{442}
=\frac{221}{442}
=\frac12.
\]

Hence the expected number of hearts is
\[
\boxed{\frac12}.
\]
\end{solution}

\begin{problem}[At least two heads]
A biased coin lands heads with probability $0.4$. It is tossed $4$ times. Find the probability of at least two heads.
\end{problem}

\begin{hint}
Add $P(X=2)$, $P(X=3)$, and $P(X=4)$, or subtract $P(X=0)$ and $P(X=1)$.
\end{hint}

\begin{solution}
Let $X\sim \mathrm{Bin}(4,0.4)$. Then
\[
P(X\ge 2)=1-[P(X=0)+P(X=1)].
\]
Now
\[
P(X=0)=(0.6)^4=0.1296,
\]
and
\[
P(X=1)=\binom41(0.4)(0.6)^3=4(0.4)(0.216)=0.3456.
\]
Hence
\[
P(X\ge 2)=1-0.1296-0.3456=0.5248.
\]
\end{solution}

\begin{problem}[Exactly one black card]
Two cards are drawn from a standard deck without replacement. Find the probability that exactly one card is black.
\end{problem}

\begin{hint}
Count one black and one red in either order, or use combinations.
\end{hint}

\begin{solution}
Using combinations,
\[
P(\text{exactly one black})=\frac{\binom{26}{1}\binom{26}{1}}{\binom{52}{2}}
=\frac{676}{1326}=\frac{26}{51}.
\]
\end{solution}

\begin{problem}[Expected number of tails]
A fair coin is tossed $9$ times. Let $X$ be the number of tails. Find $E(X)$.
\end{problem}

\begin{hint}
Recognise a binomial model.
\end{hint}

\begin{solution}
\[
X\sim \mathrm{Bin}\left(9,\frac12\right),
\]
so
\[
E(X)=np=9\cdot \frac12=\frac92=4.5.
\]
\end{solution}

\begin{problem}[Fair coin tossed $n$ times]
If a fair coin is tossed $n$ times, where $n>1$, find the probability of obtaining:
\begin{enumerate}[(i)]
  \item $n$ tails,
  \item at least two tails,
  \item no more than half of the tosses as tails.
\end{enumerate}
\end{problem}

\begin{hint}
Let $X$ be the number of tails, so $X\sim \mathrm{Bin}\left(n,\frac12\right)$. For part (iii), use symmetry about the middle and treat odd and even $n$ separately.
\end{hint}

\begin{solution}
Let
\[
X\sim \mathrm{Bin}\left(n,\frac12\right),
\]
where $X$ is the number of tails.

\textbf{(i)} Getting $n$ tails means every toss is a tail, so
\[
P(X=n)=\left(\frac12\right)^n.
\]

\textbf{(ii)} ``At least two tails'' means $X\ge 2$, so
\[
P(X\ge 2)=1-[P(X=0)+P(X=1)].
\]
Now
\[
P(X=0)=\left(\frac12\right)^n
\]
and
\[
P(X=1)=\binom n1 \left(\frac12\right)^n = \frac{n}{2^n}.
\]
Hence
\[
P(X\ge 2)=1-\frac{1+n}{2^n}.
\]

\textbf{(iii)} We want
\[
P\left(X\le \frac n2\right).
\]
Because the coin is fair, the binomial distribution is symmetric.

If $n$ is odd, there is no middle value, so exactly half the outcomes lie on each side:
\[
P\left(X\le \frac n2\right)=\frac12.
\]

If $n$ is even, the middle value $X=\frac n2$ must be shared equally between the two sides, so
\[
P\left(X\le \frac n2\right)=\frac12+\frac12 P\left(X=\frac n2\right).
\]
Therefore
\[
P\left(X\le \frac n2\right)=\frac12+\frac12\binom{n}{n/2}\left(\frac12\right)^n
=\frac12+\frac{1}{2^{n+1}}\binom{n}{n/2}.
\]
\end{solution}

\begin{problem}[HSC 2024-style: Driver test]
To pass a driver knowledge test, a learner must answer at least $9$ out of $10$ questions correctly.

A learner already knows the answers to $5$ questions and guesses randomly on the remaining $5$ questions. Each guessed question has $4$ options, with exactly one correct answer.

Find the probability that the learner passes the test.
\end{problem}

\begin{hint}
The learner already has $5$ marks, so passing means getting at least $4$ of the $5$ guessed questions correct.
\end{hint}

\begin{solution}
Let $X$ be the number of guessed questions answered correctly. Then
\[
X\sim \mathrm{Bin}\left(5,\frac14\right).
\]

To pass, the learner needs at least $4$ correct guesses:
\[
P(\text{pass})=P(X\ge 4)=P(X=4)+P(X=5).
\]

\[
P(X=4)=\binom54\left(\frac14\right)^4\left(\frac34\right)
=5\cdot \frac1{256}\cdot \frac34=\frac{15}{1024},
\]
and
\[
P(X=5)=\left(\frac14\right)^5=\frac1{1024}.
\]

So
\[
P(\text{pass})=\frac{15}{1024}+\frac1{1024}=\frac{16}{1024}=\frac1{64}.
\]
\end{solution}

\begin{problem}[HSC 2024-style: Same-metal coins]
A bag contains $7$ silver coins and $3$ bronze coins. Two coins are drawn without replacement. Find the probability that both coins are made of the same metal.
\end{problem}

\begin{hint}
Add the probabilities of both silver and both bronze.
\end{hint}

\begin{solution}
\[
P(\text{same metal})=P(\text{both silver})+P(\text{both bronze}).
\]

\[
P(\text{both silver})=\frac7{10}\cdot\frac69=\frac{42}{90},
\qquad
P(\text{both bronze})=\frac3{10}\cdot\frac29=\frac6{90}.
\]

Therefore
\[
P(\text{same metal})=\frac{42}{90}+\frac6{90}=\frac{48}{90}=\frac{8}{15}.
\]
\end{solution}

\begin{problem}[HSC 2023-style: Mixed equipment probabilities]
A gym has $5$ treadmills and $4$ rowing machines. On average, each treadmill is in use with probability $0.65$ and each rowing machine is in use with probability $0.40$, independently of the others.
\begin{enumerate}[(a)]
  \item Find an expression for the probability that exactly $3$ treadmills are in use.
  \item Find an expression for the probability that exactly $3$ treadmills are in use and no rowing machines are in use.
\end{enumerate}
\end{problem}

\begin{hint}
Use a binomial expression for the treadmills, then multiply by the rowing-machine condition.
\end{hint}

\begin{solution}
\textbf{(a)} For the $5$ treadmills,
\[
P(\text{exactly 3 in use})=\binom53(0.65)^3(0.35)^2.
\]

\textbf{(b)} The probability that no rowing machines are in use is
\[
(0.60)^4.
\]
So the required expression is
\[
\binom53(0.65)^3(0.35)^2(0.60)^4.
\]
\end{solution}

\begin{problem}[HSC 2021-style: A Bernoulli complement]
The random variable $X$ represents the number of successes in $10$ independent Bernoulli trials, where the probability of success in each trial is $0.9$.

If
\[
r=P(X\ge 1),
\]
determine whether $r$ is greater than, equal to, or less than $0.9$.
\end{problem}

\begin{hint}
Use the complement $P(X\ge 1)=1-P(X=0)$.
\end{hint}

\begin{solution}
\[
r=P(X\ge 1)=1-P(X=0)=1-(0.1)^{10}.
\]
Since $(0.1)^{10}$ is extremely small,
\[
r=1-(0.1)^{10}>0.9.
\]
\end{solution}
